\section{Methods}
\label{sec:methods}

This section describes the experimental design, model architectures, projection
optimizer variants, and the two primary engineering contributions of this work:
projection-scoped normalization and the \texttt{iFOPNG} algorithm.

\subsection{Benchmarks and Data}
All experiments use two standard sequential classification benchmarks.
\textbf{Split-MNIST} divides the 10-digit MNIST dataset
($28{\times}28$ grayscale; $60{,}000$ training samples) into five binary
classification tasks—digits $\{0,1\}$, $\{2,3\}$, $\{4,5\}$, $\{6,7\}$, and
$\{8,9\}$—presented in sequence.  \textbf{Split-CIFAR10} divides the 10-class
CIFAR-10 dataset ($32{\times}32$ RGB; $50{,}000$ training samples) into five
binary tasks in the same fashion.  Both streams operate under the
\emph{incremental task learning} protocol with known task identity at test time:
each task receives a dedicated output head, and the model is evaluated on
all previously seen tasks after each new task is trained.

\subsection{Model Architectures}
\paragraph{Standalone Multi-Head CNN.}
As an unconstrained upper-bound reference, a vanilla multi-head convolutional
network is trained directly on the target task stream.  This architecture
serves as a ceiling against which hypernetwork-based approaches are compared.

\paragraph{Chunked Hypernetwork.}
The hypernetwork $h(\cdot;\,\phi)$ generates the complete weight vector
$\theta_t = h(e_t;\,\phi)$ of a single-head target CNN conditioned on a
task embedding $e_t \in \mathbb{R}^{d_t}$.  Because generating the full
$\theta_t$ in a single forward pass would require an output layer of dimension
$|\theta_t|$ (tens of thousands of scalars), the hypernetwork instead employs
\emph{weight chunking}: a secondary chunk embedding
$c_k \in \mathbb{R}^{d_c}$, $k=1,\ldots,K$, is concatenated to $e_t$, and the
network produces a block of $C_s$ parameters per forward pass.  The complete
target weight is assembled by iterating over all $K$ chunks.  The functional
regularizer of \citet{oswald_continual_2022} is applied on the generated
weights:
\begin{equation}
    L_{\text{total}} = L_t(\theta_t)
    + \beta \sum_{s < t} \bigl\|h(e_s;\,\phi) - \theta_s^*\bigr\|_2^2.
\end{equation}
Unless stated otherwise, all hypernetwork experiments use the same shared
architecture: hidden dimension 32, chunk size $C_s = 64$,
$d_t = d_c = 32$, trained for up to 50 epochs per task.

\subsection{Projection Optimizer Suite}
Six projection optimizer variants are compared alongside the baselines.
\textbf{OGD} \citep{farajtabar_orthogonal_2019} stores raw gradient directions
and projects via Euclidean orthogonality.
\textbf{FNG} uses the Fisher metric for preconditioning without an orthogonal
projection constraint.
\textbf{ONG} combines natural gradient descent with Euclidean orthogonal
projection, analogously to FOPNG but without the Fisher inner product in the
projection kernel.
\textbf{FOPNG} \citep{garg_fisher-orthogonal_2026} executes the full
Fisher-orthogonal projection:
\begin{equation}
    \Delta\theta = -\eta
    \bigl(I - G(G^\top F G)^{-1} G^\top F\bigr)
    F^{-1}\nabla_\theta L_t,
\end{equation}
where $G = [g_1 \mid \cdots \mid g_{t-1}]$ is the memory matrix of historical
gradient directions and $F$ is the diagonal empirical Fisher.
\textbf{iFOPNG} (ours) modifies the preconditioner; its derivation is presented
in Section~\ref{sec:ifopng}.
The two baselines are \textbf{Adam} (Adam with $L_2$ weight decay folded into
the gradient) and \textbf{SGD} with momentum, both tuned with learning rate
$\eta = 10^{-3}$; all projection optimizers use $\eta = 10^{-2}$, $K = 250$
stored gradients per task, and up to 50 training epochs.  See Appendix~B for
the complete hyperparameter table.

\subsection{Projection-Scoped Normalization (Sub-RQ2)}
\label{sec:normalization}

The chunk-induced gradient pathology (Section~\ref{sec:related_work}) arises
because the hypernetwork executes $K$ forward–backward passes before a single
projection step, causing the gradient columns of $G$ and the diagonal Fisher
$\hat{F}$ to accumulate to extreme magnitudes.  When these unnormalized vectors
are used in the projection kernel $(G^\top F G)^{-1}$, the resulting condition
number exceeds machine precision, causing inversion to fail catastrophically.

We resolve this by applying two local rescaling operations \emph{within} the
projection computation and nowhere else, ensuring the underlying optimization
dynamics are not altered:

\begin{enumerate}
    \item \textbf{Gradient column normalization.}  Each new gradient vector
    $g_t$ is rescaled to unit $\ell_2$ norm before insertion into $G$:
    \begin{equation}
        \tilde{g}_t = \frac{g_t}{\|g_t\|_2 + \varepsilon}, \quad \varepsilon = 10^{-8}.
    \end{equation}
    This ensures that all historical directions contribute equally to the span
    of $G$, independent of the magnitude explosion induced by chunking.
    Any SVD-based compression of $G$ is subsequently applied to these unit
    vectors, so the subspace geometry is preserved exactly.

    \item \textbf{Fisher absolute-maximum normalization.}  The diagonal
    empirical Fisher $\hat{F}$ is rescaled by its own maximum entry:
    \begin{equation}
        \tilde{F} = \frac{\hat{F}}{\max_i |\hat{F}_i| + \varepsilon}.
    \end{equation}
    This keeps the Fisher entries in $[0, 1]$, ensuring that the inner-product
    form $G^\top F G$ has a well-conditioned scale regardless of the
    number of chunks.
\end{enumerate}
These two normalizations together prevent the condition number blow-up,
completely resolving inversion failures in all experiments.

\subsection{First-Task Initialization: AdamW versus Adam (Sub-RQ3)}
\label{sec:adamw}

Because the projection memory matrix $G$ is populated exclusively during the
first task, the quality of the initial subspace is determined by the optimizer
used at that stage.  Standard \texttt{Adam} with $L_2$ weight decay folds the
decay term \emph{into the gradient estimate}, meaning the historical gradient
vectors stored in $G$ are contaminated by a component proportional to the
current weight magnitudes rather than the task loss curvature alone.  This
contaminates the projection basis with a weight-magnitude artifact unrelated to
task geometry.

\texttt{AdamW} \citep{loshchilov_decoupled_2019} decouples weight decay from
the gradient update, applying it as a direct parameter shrinkage step after the
gradient-based update.  Consequently, the gradients stored in $G$ reflect only
the task loss curvature, keeping the projection subspace geometrically clean.
All projection optimizer runs in this work use \texttt{AdamW} for the first task
and the respective projection algorithm for subsequent tasks.

\subsection{iFOPNG: Elastic Fisher-Orthogonal Projected Natural Gradient}
\label{sec:ifopng}

Rigid FOPNG enforces a hard boundary defined by $F_{\text{new}}$, the
Fisher information of the current task alone.  Parameters that were important
in old tasks but carry low current-task Fisher mass are therefore free to drift
once a few tasks have passed, which may erode the protection conferred by the
functional regularizer.

To address this, \texttt{iFOPNG} replaces $F_{\text{new}}$ with the combined
curvature $F_c = F_{\text{new}} + F_{\text{old}}$, where $F_{\text{old}}$ is an
accumulated diagonal Fisher across all prior tasks, maintained as the element-wise
maximum of the historical per-task Fisher vectors.  The iFOPNG update solves the
constrained optimization:
\begin{align}
    \min_{u \in \mathbb{R}^p} \quad
        & \| g - F_c^{-1} u \|_{F_c}^2 \\
    \text{s.t.} \quad
        & F_{\text{old}}^{-1} u \in \operatorname{Col}(G),
\end{align}
yielding the closed-form parameter update:
\begin{equation}
    v^* = \eta\,
        \frac{F_c^{-1} P g}{\sqrt{(Pg)^\top F_c^{-1} (Pg)}},
    \qquad
    P = I - F_{\text{old}}\,G\,A^{-1}\,G^\top F_{\text{old}},
    \label{eq:ifopng}
\end{equation}
where $A = G^\top F_{\text{old}}\,F_c^{-1}\,F_{\text{old}}\,G$.  A formal
derivation is provided in Appendix~A.  Intuitively, $F_c$ accumulates curvature
from all tasks seen so far: parameters important to old tasks remain damped in
the natural-gradient trust region even when $F_{\text{new}}$ alone would not
penalize moving them, providing an EWC-style elastic inertia boundary without
requiring an explicit regularization term.

\subsection{Evaluation Metrics and Implementation}
Model performance is quantified via two metrics computed at the end of each task:
(i)~\textbf{Average accuracy} over all tasks trained so far, and
(ii)~\textbf{Backward Transfer (BWT)}, defined as the mean accuracy change on
previously trained tasks:
\begin{equation}
    \mathrm{BWT} = \frac{1}{T-1}\sum_{t=2}^{T}
    \bigl(\text{acc}_{T,t} - \text{acc}_{t,t}\bigr),
\end{equation}
where $\text{acc}_{i,j}$ is the accuracy on task $j$ evaluated after training
task $i$.  All experiments are run on CUDA-enabled hardware.  Large projection
matrices and Fisher vectors are offloaded to CPU in a hybrid memory management
scheme to support the high gradient budgets required ($K = 250$ gradients per
task) within realistic VRAM constraints.

% ─────────────────────────────────────────────────────────────────────────────